\section{Laminar Chemistry Study\label{ch4:laminarChemistryStudy}}

\todo{Use Results from AIAAA, laminar comparison between frozen and
nonequilibrium flow}
This section investigates nonequilibrium effects of aero-optics in a laminar
flow by comparing a flow with nonequilibrium chemistry and the other one with
frozen chemistry; \TabRef{tab:laminarChemistryStudy} shows the cases used for this
investigation. 

\begin{table}[H]
    \begin{center}
        \begin{tabular}{|c|c|c|c|c|c|c|}
            \hline
            Case & $M\Units{\;}$ & $T_{wall}\Units{K}$ &
            $T_{\infty}\Units{K}$ & $P_{\infty}\Units{Pa}$ &
            $h_0\Units{MJ}$ & $q\Units{kPa}$ \\
            \hline\hline
            $1C$\cite{Mackey2019} & $11.2$ & $1000$ & $293$ & $1200$ & $7.38$  & $105$ \\\hline
            $2C$\cite{Mackey2019} & $13.2$ & $1000$ & $293$ & $1200$ & $10.25$ & $146$ \\\hline
            $3C$\cite{Mackey2019} & $15.2$ & $1000$ & $293$ & $1200$ & $13.60$ & $193$ \\\hline
        \end{tabular}
    \caption{Chemistry Study \label{tab:laminarChemistryStudy}}
    \end{center}
\end{table}

    For this investigation, the IRV-2 vehicle, shown in \FigRef{fig:IRV2} is used.
    The IRV-2 vehicle comprises a sphere-biocone shape with a nose radius of $1.95\Units{cm}$ and a total
    length of $1.386 \Units{m}$ \cite{Kuntz2001}. 

    \pic{0.5}{chapter5/laminarChemistry/irv2}{IRV-2 vehicle}{fig:IRV2}

\todo{Discuss about IRV and mesh}

